% New English translation from checked cote 25, site 81 (PDF page 82).
% AI-translated by GPT-6 Astra (OpenAI, ChatGPT Pro mode).
\scanpage{25/81}
\mspage{\hand{76}}

NB \quad The argument no longer works if we replace x by a
linear subvariety \unsure{$X$|$C$} of dim.$\hand{>}0$, and require H to pass
through C: indeed, nothing prevents (taking for example dim C = r-2, the
largest possible so that H can still actually vary)
T from containing the line $C^{o}$:\var{\xout{(and to give an example where where this occurs}}{}{typewriter strike-out}{ega25_25_081-v0001}
\var{\xout{in char 0, say, we start from a subvariety T of}}{}{typewriter strike-out}{ega25_25_081-v0002} take, for
example, a \var{\inserted[typed]{nonsingular}}{nonsingular}{insertion}{ega25_25_081-v0003} quadric in $P$, in arbitrary characteristic. But it
is possible that such phenomena cannot occur either \var{\xout{\unsure{if}}}{}{typewriter strike-out}{ega25_25_081-v0004}
for $\hand{\phi}_n f$, $n\hand{\geq} 2$; we could pose the question in a remark in
no. 9.

\uline{Remarks} 13.4. \quad a) We have already pointed out that the hypothesis that k \var{\xout{\illegible{1}}}{}{typewriter strike-out}{ega25_25_081-v0005}
is perfect is essential for the validity of 13.2. On the other hand, it
is plausible that the hypothesis “k infinite” is superfluous. \var{\xout{In any case}}{}{typewriter strike-out}{ega25_25_081-v0006}
We have not tried \var{\xout{here}}{}{typewriter strike-out}{ega25_25_081-v0007} to give an ad hoc argument for the case
where k is finite, and merely point out that in this case, applying
13.2 to the algebraic closure of k and using the usual arguments
shows that we can find a finite extension k' of k such that
\var{\xout{\unsure{at}}}{}{typewriter strike-out}{ega25_25_081-v0008} \var{\inserted[typed]{for}}{for}{insertion}{ega25_25_081-v0009} the points of $X_{k'}$ over x there exist open neighborhoods
that are poly-elementary relative to k'.

b) When in 13.2 we drop the hypothesis that X is geom
irreducible, the conclusion of course no longer holds as stated
(since a poly-elementary algebraic scheme is geom irred!). It
does however hold in the weaker form that we would obtain by
omitting the word “connected” from definition 13.1: this is what
the proof we have given shows.
\var{\inserted{c)}}{c)}{insertion}{ega25_25_081-v0010}

\var{\over{c}{d}}{d}{writing over}{ega25_25_081-v0011}) Proceeding as explained in 7.1, we can \var{\inserted{give}}{give}{insertion}{ega25_25_081-v0012} variants
of 13.2 in the case where we replace the base field k by a prescheme
\var{\margin{81}}{\placedmargin{81}}{marginal note placed}{ega25_25_081-v0013}
\tnote{25-081-1}{The opening X/C alternatives, the repeated where in the canceled sentence, and the empty inserted item c) before the correction of c) to d) are retained.}
